%% kapa: exam=1 round=37 year=2026 period=1 form=A subject=civil_law subject_ko=민법 seq=19 number=19
%% source: Q-Net 감정평가사 기출문제 (cst00309) artlSeq=5258924 file=raw/1008/5258924/2026년 제37회 감정평가사 1차 1교시.pdf
%% conversion: kordoc 4.24.51 -> tools/gen1.py (auto); review=proofread-sonnet; flags=none
\begin{question}{19}
소멸시효에 관한 설명으로 옳지 않은 것은? (다툼이 있으면 판례에 따름)
\begin{choices}
  \item 세무사의 직무에 관한 채권은 10년의 소멸시효에 걸린다.
  \item 채무자가 시효완성 후 채무를 승인한 경우에는 시효완성 사실을 알고 그 이익을 포기한 것으로 보아야 한다.
  \item 판결에 의하여 확정된 채권은 단기의 소멸시효에 해당한 것이라도 그 소멸시효는 10년으로 한다.
  \item 부작위를 목적으로 하는 채권의 소멸시효는 위반행위를 한 때로부터 진행한다.
  \item 후순위 담보권자는 선순위 담보권의 피담보채권 소멸로 직접 이익을 받는 자에 해당하지 않으므로 선순위 담보권의 피담보채권에 관한 소멸시효가 완성되었다고 주장할 수 없다.
\end{choices}
\end{question}
\answer{2}
